% SPDX-License-Identifier: CC-BY-4.0
\section{Experimental Setup}
\label{sec:setup}

\subsection{Two Corpora}
\label{sec:corpora}

We evaluate every system on two corpora.

\textbf{Own corpus (N=34).} Curated for the project, organised by
intent and schema (4 dialogue queries with cascade-derived seeds,
2 grounded-assertion queries, 2 empty-working-set queries, 2
high-confabulation-risk queries, 1 shallow-cascade query, 1
valence-opposed query, plus three fluency-exercise queries from
ADR-0008, ten bucket-coverage queries from ADR-0009 across
Conversational/Formal/Neutral/Mixed registers, and nine
schema-coverage queries from ADRs-0016, 0023, 0024, 0025 across
Summary, ComparativeAnalysis, IntrospectiveAssessment, and
NarrativeArc). The corpus is concentrated on engineering and
deployment scenarios reflecting the project's primary intended
deployment context. Source:
\texttt{crates/hyphae-eval/src/corpus.rs}.

\textbf{TriviaQA subset (N=150).} Random sample (seed 42) from
the TriviaQA \texttt{rc} validation split
\citep{joshi2017triviaqa}. For each sampled query, we extract a
seed body from the supplied Wikipedia context by (1) filtering to
sentences containing the answer or any alias as a word-bounded
match, (2) restricting to sentences of 30--250 characters, and
(3) reranking the survivors by cosine similarity to the query
using the \texttt{sentence-transformers/all-MiniLM-L6-v2}
embedder. The highest-scoring sentence becomes the seed body. The
filter rejection rate on this seed is approximately 17\%
(21 samples without wiki\_context, 10 without a qualifying
sentence) yielding 150 surviving queries. Source:
\texttt{bench/baseline-llm-rag/src/baseline\_llm\_rag/corpus\_external.py}.

The two corpora are complementary. The own corpus tests
multi-fragment composition (working sets of 2--3 seeds with
varied registers); TriviaQA tests single-fragment retrieval. The
own corpus is dense with engineering scenarios; TriviaQA spans
biography, geography, history, science, pop culture, and sports.

\subsection{Generator Models}
\label{sec:generators}

We compare \hyphae{} against six LLM-based configurations
spanning the model-class landscape:

\begin{itemize}
    \item \textbf{Llama-3.1-8B-Instruct Q4\_K\_M} (4.6 GB GGUF):
    small open model, run locally via
    \texttt{llama-cpp-python} on the laptop's Metal backend or a
    DigitalOcean CPU droplet. Establishes the small-open
    baseline.
    \item \textbf{Llama-3.3-70B-Instruct} (FP8 hosted): larger
    open model, accessed via DigitalOcean Inference
    \citep{do_inference}. Tests whether scale helps.
    \item \textbf{Claude-4.6-Sonnet} (frontier closed,
    Anthropic): tests a frontier model from a different lab
    family.
    \item \textbf{GPT-4.1} (frontier closed, OpenAI): tests a
    frontier model from the family most-frequently cited as the
    LLM baseline.
    \item \textbf{DeepSeek-V4-Pro} (frontier open reasoning):
    tests a frontier open model with reasoning training.
    \item \textbf{router:celiums-conversation} (our in-house
    Atlas router): tests a production LLM-routing system
    configured for conversational dispatch.
\end{itemize}

All six are accessed through a uniform OpenAI-compatible API
endpoint provided by DigitalOcean's GenAI Platform, with the
exception of the local Llama-8B which uses \texttt{llama.cpp}.
This uniformity removes API-shape variance from the comparison
while preserving the providers' native decoding pipelines.

\subsection{Retrieval Modes}

Each LLM is evaluated under three retrieval modes:

\begin{itemize}
    \item \textbf{Oracle.} The corpus's seed bodies are passed
    directly as context. No retrieval; the LLM receives the gold
    supporting passage. Isolates composition quality from
    retrieval quality.
    \item \textbf{Vanilla RAG.} FAISS \citep{johnson2019faiss}
    \texttt{IndexFlatIP} (exact cosine after L2 normalisation)
    over chunks of all seed bodies. Top-$k=5$ retrieved, passed
    as context. Chunking uses \texttt{tiktoken cl100k\_base}
    BPE with 256-token chunks and 32-token overlap.
    \item \textbf{Strong RAG.} HyDE plus cross-encoder reranking.
    The LLM first generates a hypothetical answer; that answer is
    embedded and retrieved (over-retrieve $k=20$); the candidates
    are reranked by the \texttt{BAAI/bge-reranker-base}
    cross-encoder \citep{bge_reranker} and the top-5 passed as
    context.
\end{itemize}

The retrieval pipeline (embedder, chunker, FAISS index, reranker)
is shared across all six LLMs. Only the generation stage varies.

\subsection{Metrics}
\label{sec:metrics}

We score every system on the same set of metrics, comprising
those directly comparable across architectures (the
\emph{comparable subset}) and two extra metrics added to capture
properties specific to the verbatim-quotation commitment.

\paragraph{Comparable subset.}

\begin{itemize}
    \item \verbpass{} (boolean): every seed body marked
    \texttt{verbatim\_quotation} appears verbatim in the
    response.
    \item \texttt{connective\_hygiene\_pass} (boolean): the
    response is free of doubled-connective stutters from a
    pinned list.
    \item \texttt{ngram\_overlap\_$n$} ($n \in \{4, 5, 8\}$): the
    fraction of $n$-grams in the response that appear in the
    concatenated retrieved chunks after lowercasing and
    whitespace normalisation.
    \item \unsupf{} (NLI-based): each response is split into
    sentences. Sentences beginning with a known connective
    phrase from a pinned list are excluded as scaffolding. The
    remaining factual sentences are scored by
    \texttt{roberta-large-mnli} as entailment / neutral /
    contradiction against the concatenated retrieved context.
    The filtered rate is (neutral + contradiction)~/~(total
    factual).
    \item \unsupr{}: the same as \unsupf{} without the connective-
    sentence filter. Higher than \unsupf{} for any system whose
    output contains scaffolding phrases.
    \item \texttt{latency\_p50}, \texttt{latency\_p95},
    \texttt{latency\_mean}: per-query wall-clock latency
    aggregated over the corpus.
\end{itemize}

\paragraph{Verbatim-quotation extras.}

\begin{itemize}
    \item \texttt{quoted\_content\_supported\_rate}: the fraction
    of double-quoted spans in the response that appear verbatim
    in the retrieved context. Hyphae produces quoted spans for
    every fragment; LLM-based systems rarely use double quotes
    formally.
\end{itemize}

\paragraph{Hyphae-only diagnostics.}

\hyphae{} produces a structured \texttt{RealizationOutput}
carrying limitation triggers, schema identifier, and
fragments-quoted list. The internal scoring harness measures
nine additional dimensions over this output (schema match,
limitation recall/precision, lexical diversity, role coverage,
boundary smoothness). These are inapplicable to LLM-based
systems and are reported separately, not in the head-to-head
table.

\paragraph{Statistical inference.}

For every aggregate metric we report a bootstrap percentile
confidence interval (1000 resamples, 95\% level, seed=42). At
$N=34$ (own corpus) the CIs are wide; at $N=150$ (TriviaQA) they
narrow substantially. The headline results we emphasise survive
this narrowing.

\subsection{Hardware}
\label{sec:hardware}

The primary experimental hardware is an Apple Silicon laptop
(M-series, 10 cores, MPS backend for the local Llama-8B and the
NLI scorer). Network calls to DigitalOcean Inference run on
the laptop's network.

For the hardware-matrix portion (Section~\ref{sec:results-hw})
we provisioned a DigitalOcean \texttt{c-16} CPU-optimised
droplet (Intel Xeon Platinum 8168, 16 dedicated vCPU, 31 GB
RAM, Ubuntu 24.04 LTS, NYC1 region) and re-ran the local Llama-8B
matrix and the \hyphae{} ablations there.

The reported \hyphae{} latencies should be interpreted as a
lower bound on what is achievable: optimisation work has not
been prioritised. The reported LLM latencies should be
interpreted as the best plausible API-routed latency under no
batching --- production systems with batching and rate-limited
parallelism may achieve different throughput, but the per-query
latency on the user-perceived path is bounded below by the
LLM's serial inference time.
